\section{Illustrative scan of deficits}
A negative auxiliary signal amplitude describes missing events with the signal-template shape; it is not a physical negative rate or coupling. This figure mirrors the excess scan without changing its fits or limits.

\fig{0.97}{../../v4p9p16_deficit_extension_20260906/figures/combined_deficit_scan.pdf}{Profiled deficit scan. Top: observed signed roots and the stress-background offset; shading marks negative fitted contributions. Middle: raw-root Gaussian reference and conditional Gaussian-response tails. Both assign one to nonnegative raw roots. Local triangles mark the $10^{-8}$ display floor. Bottom: two separate global orderings; triangles mark one-sided 95\% zero-count bounds. Direct error bars are central 95\% intervals.}{fig:deficit}

For $z=(r-a)/s$, the conditional local rule is $p^-=\Phi(z)$ when $r<0$, and one otherwise. The principal scan score is $T^-=\max_{m:r_m<0}(-z_m)$, with $-\infty$ for an empty set. The separate raw-depth statistic is $D^-=\max_m\max(0,-r_m)$.

The deepest raw profiled deficit is at 72 MeV, with $r=-3.490$ and an uncalibrated local reference of $2.41\times10^{-4}$. The largest stress-centered deficit instead occurs at 83 MeV: $r=-0.676$, $a=+7.707$ and $s=0.983$. Its raw reference is 0.2495, whereas the conditional local value is $7.18\times10^{-18}$. This contrast again comes mainly from the stress offset.

The latter threshold has 0/200,000 GP and 0/1,000 direct exceedances. The one-sided 95\% global upper bounds are $1.50\times10^{-5}$ and 0.003, respectively. Every simulated raw-depth maximum exceeds the observed deficit. These are the same realizations used for the excess scan; no new fits or independent toys were needed.

These illustrative tails follow the excess scan and do not adjust for choosing directions, methods or orderings. The unresolved tails establish neither a particle signal nor physical background validity.
